% operator-stability-v1 — publication template (writeup-aligned)
% Slots filled by scripts/generate_operator_paper.py
\documentclass[12pt]{article}
\usepackage[letterpaper,top=2cm,bottom=2cm,left=3cm,right=3cm]{geometry}
\usepackage[T1]{fontenc}
\usepackage[utf8]{inputenc}
\usepackage{lmodern}
\usepackage{microtype}
\usepackage{amsmath,amssymb,amsthm,mathtools}
\usepackage{booktabs}
\usepackage{xcolor}
\usepackage{enumitem}
\usepackage{xurl}
\usepackage[colorlinks=true,linkcolor=blue,urlcolor=blue,citecolor=blue]{hyperref}

\Urlmuskip=0mu plus 1mu
\urlstyle{tt}
\emergencystretch=2em

\setlength{\parskip}{0.4em}
\setlength{\parindent}{1.2em}
\setlist[itemize]{leftmargin=1.4em,itemsep=0.25em,topsep=0.35em}
\setlist[enumerate]{leftmargin=1.4em,itemsep=0.25em,topsep=0.35em}

\theoremstyle{plain}
\newtheorem{theorem}{Theorem}
\newtheorem{lemma}[theorem]{Lemma}
\theoremstyle{definition}
\newtheorem{definition}[theorem]{Definition}
\newtheorem{example}[theorem]{Example}
\theoremstyle{remark}
\newtheorem{remark}[theorem]{Remark}

% Breakable monospace for Lean paths / identifiers (xurl \path)
\newcommand{\leanpath}[1]{%
  {\raggedright\urlstyle{tt}\path{#1}}}


\definecolor{fundbox}{RGB}{245,245,248}
\definecolor{fundrule}{RGB}{180,180,190}

\title{Constraint-threshold conjunction: Real Stability}
\author{Research Architecture Runtime\\[0.25em]
{\normalsize\texttt{operator/constraint-threshold-conjunction}}}
\date{2026-07-29}

\begin{document}
\maketitle

\begin{abstract}
\leanpath{constraint-threshold-conjunction} preserves a Boolean combination of threshold constraints under buffered scalar noise. Lean \texttt{LEAN\_FULL}.
\end{abstract}

\noindent
\begin{center}
\fcolorbox{fundrule}{fundbox}{%
  \begin{minipage}{0.92\linewidth}
  \centering
  \textbf{This theorem is:} \texttt{Reduction}
  \end{minipage}}
\end{center}
\vspace{0.6em}

\section{Problem}
{\raggedright\sloppy
Conjunction or disjunction of threshold predicates on coordinates/scores.
\par}

\section{Stability notion}
Each literal is threshold-stable; Boolean combinations inherit preservation.

\section{Definitions}
Reduce to multi-threshold / threshold preservation on each cut.

\section{Theorem}
\begin{theorem}\label{thm:inv}
If every atomic threshold is buffer-stable for $(x,x',\varepsilon)$, the conjunction/disjunction value is preserved.
\end{theorem}

\section{Intuition}
Boolean operations are deterministic functions of stable bits.

\section{Examples}
\begin{example}
Two cuts $x\ge 1$ and $x\ge 3$ with $x=4$, $\varepsilon=0.2$: both bits stable, so their conjunction stays true.
\end{example}

\section{Proof}
Each atomic predicate is preserved by scalar threshold preservation under the buffer hypotheses. Conjunction (resp.\ disjunction) is a fixed Boolean function of those bits, hence preserved (Theorem~\ref{thm:inv}).

\section{Formal statement}
{\raggedright\sloppy
\begin{description}[style=nextline,leftmargin=1.1em,font=\normalfont\bfseries]
\item[Lean module] \leanpath{Research.Operators.ConstraintThresholdConjunction.Preservation}
\item[Source file] \leanpath{lean/Research/Operators/ConstraintThresholdConjunction/Preservation.lean}
\item[Theorems]\leavevmode\par\vspace{0.15em}\begin{itemize}[leftmargin=1.2em]
  \item \leanpath{constraint\_threshold\_conjunction\_conjunction\_preservation}
  \item \leanpath{constraint\_threshold\_conjunction\_conjunction\_sharpness}
\end{itemize}
\item[Certificate] \leanpath{lean/certificates/constraint-threshold-conjunction/constraint-threshold-conjunction-conjunction-preservation}
\item[Status] \texttt{LEAN\_FULL} on Mathlib $\mathbb{R}$ (\texttt{REAL\_MATHLIB})
\end{description}
\par}

\section{Proof dependencies}
{\raggedright\sloppy
\begin{itemize}\item Threshold / multi-threshold cores.\item Boolean congruence.\end{itemize}
\par}

\section{Consequences}
{\raggedright\sloppy
Constraint-side companion to multi-threshold counts.
\par}

\end{document}
